\chapter{Risk Assessment Report}
\label{chap:risk_assessment_report}

This chapter summarises the output of the formal risk assessment carried out for GreenGrid Energy.  
The process complies with ISO/IEC 27001 Clause 8.2 and uses the methodology specified in Chapter~\ref{chap:risk_management}, applied to the critical assets listed in Appendix~\ref{appendix:asset_inventory}.  
Likelihood and impact were estimated with the GreenGrid Energy risk-matrix shown in Figure~\ref{figure:risk_matrix}; the resulting risk levels drive the Statement-of-Applicability in Chapter~\ref{chap:soa}.

%-----------------------------------------------------------------------
\section{Risk-assessment results}

The subsections below group risks by \emph{shared} assets—resources consumed by several departments—and by assets that are \emph{unique} to Supply-Chain Management (SCM).  
This separation makes residual-risk ownership clear and avoids duplicated treatment plans.

%-----------------------------------------------------------------------
\subsection{Shared assets}

Shared assets include the company-wide ERP instance, the central IT network stack and services managed by HR and Customer Services.  
Because these systems cut across functional boundaries, a successful attack could propagate rapidly, amplifying both impact and remediation cost.  
Table~\ref{tab:shared_risk} lists the highest-ranked risks for those assets.

\begin{longtable}{|p{2cm}|p{2.5cm}|p{3cm}|p{2cm}|p{2cm}|p{2cm}|}
  \hline
  \textbf{Asset ID} & \textbf{Threat} & \textbf{Vulnerability} & \textbf{Likelihood} & \textbf{Impact} & \textbf{Risk Level} \\
  \hline
  \endhead
  ERP-301 & Data Breach & Poor access controls & Possible & Severe & \textbf{Medium-High} \\
  \hline
  IT-201 & Phishing Attacks & Insufficient staff training & Very Likely & Significant & \textbf{High} \\
  \hline
  IT-202 & Unauthorised Access & Firewall misconfiguration & Possible & Significant & \textbf{Medium-High} \\
  \hline
  CS-401 & Data Leakage & Improper permission management & Possible & Significant & \textbf{Medium-High} \\
  \hline
  HR-501 & Insider Threat & Excessive user privileges & Possible & Significant & \textbf{Medium-High} \\
  \hline
  HR-502 & Staff Unavailability & High dependency on key staff & Possible & Significant & \textbf{Med-High} \\
  \hline
\end{longtable}

The assessment highlights two systemic drivers:

\begin{itemize}
  \item \textbf{Credential misuse.}  Four of the six entries trace back to weak or excessive privileges; tightening identity-and-access management (IAM) therefore offers a high return on investment.
  \item \textbf{Human factors.}  The \texttt{IT-201} phishing scenario scores \emph{High} even though no technical vulnerability is exploited—demonstrating that staff security culture is as critical as tooling.
\end{itemize}

Mitigation must therefore balance technical hardening with sustained awareness training and succession planning for key roles.

%-----------------------------------------------------------------------
\subsection{Supply-chain assets}

SCM relies almost entirely on the corporate ERP platform (\texttt{ERP-301}) to manage purchase orders, vendor master data and inbound-logistics scheduling.  
Although the asset is physically shared, its SCM data flows, integrations and uptime requirements differ from finance or HR, warranting a dedicated risk view (Table~\ref{tab:scm_risk}).

\begin{longtable}{|p{2cm}|p{2.5cm}|p{3cm}|p{2cm}|p{2cm}|p{2cm}|}
  \hline
  \textbf{Asset ID} & \textbf{Threat} & \textbf{Vulnerability} & \textbf{Likelihood} & \textbf{Impact} & \textbf{Risk Level} \\
  \hline
  \endhead
  ERP-301 & Supply Chain Compromise & Insecure third-party integrations & Possible & Severe & \textbf{Medium-High} \\
  \hline
  ERP-301 & Data Integrity Loss & Configuration errors, incorrect updates & Possible & Significant & \textbf{Medium-High} \\
  \hline
  ERP-301 & System Downtime & Single point of failure & Unlikely & Severe & \textbf{Medium-High} \\
  \hline
  \end{longtable}

Here the prime issue is \emph{external trust}. Unvetted vendor APIs or tampered firmware could inject malicious code directly into production—circumventing perimeter defences.  
Equally, a prolonged ERP outage would freeze purchase orders and delay grid-maintenance spares, exposing the organisation to regulatory penalties.  
Supplier-security clauses, robust end-to-end change management and geographically redundant hosting are therefore priority controls.

